\section{模型评价}

\subsection{模型的优点}

\begin{enumerate}
\item 两光束模型通过条纹幅度与条纹位置求解折射率与厚度信息，两条求解路径分离，每一步都可以单独检验。
\item 三个问题在同一框架下递进，两光束情形是多光束反射率式(31)的最低阶近似，多光束判别与厚度计算又各有多种独立证据相互印证。
\end{enumerate}

\subsection{模型的缺点}

\begin{enumerate}
\item 偏振按算术平均近似处理，入射角较大时存在系统偏差。
\item 峰谷法对噪声敏感，个别峰谷对的厚度偏差可达一成以上；且常数折射率在低波数端不再适用。
\item 多光束干涉下折射率只能取文献名义值，判别阈值与相干长度也依赖数值模拟与采样估算，换用其他材料与仪器时需重新标定。
\end{enumerate}

\subsection{模型的改进方向}

\begin{enumerate}
\item 用计入两种偏振分量的严格反射率模型替代算术平均近似，减弱大入射角下的近似误差。
\item 低波数段改用更精细的色散模型描述衬底折射率与吸收随波数的变化，提高深谷处的拟合精度，并使该谱段重新参与厚度统计。
\item 在多光束影响可以忽略的波段，用包络线法反演折射率，与文献名义值交叉检验。
\end{enumerate}
